\section{二维数值计算与复现方法}
\label{app:validation}
\subsection{事件树计算与条件高斯积分}
主实验枚举 $\gamma_{0:k+1}$ 的全部 64 个事件串。
对于两拍策略，先分别递推私有符号取正、负时的概率和偏差，再对同一事件串的质量求和。
因此，TV 根据符号混合后的可见记录分布计算。
本节两拍策略的条件偏差仅依赖事件串和符号，持续策略仅依赖事件历史，
故两者的总体期望均可由有限事件树精确表达；表中数值是这些表达式的双精度计算结果。
协方差直接由矩阵公式递推，不按浮点距离合并标签。

实现中分别计算 $P-P(P+R_i)^{-1}P$ 与 $(P^{-1}+R_i^{-1})^{-1}$，
以检查矩阵更新的一致性，并核对概率总和、协方差正定性、逐路径峰值、见证恒等式和有符号证书。
其中，见证恒等式通过两种独立汇总方式核对：一是对目标偏差的条件指数项求平均，
二是根据完整事件串的质量汇总下一时刻的静默位及其权重。
两种结果的一致性用于检查数值实现。

在预测均值平移为零的坐标下，另设 $x\sim N(d,P_0)$、$d=(0.13,-0.17)^\top$。
积分测量噪声后，静默似然为正文所用的高斯因子。
对 $x$ 进行二维 Gauss--Hermite 积分，独立计算静默概率、条件均值和协方差。
静默概率的解析值为 $s_{\rm cur}(P_0)\allowbreak e^{-d^\top(Z^{-1}+P_0+R)^{-1}d/2}$；
条件均值和协方差的解析值分别为 $(I-K_0(P_0))d$ 和 $F_0(P_0)$。
这里 $s_{\rm cur}(P_0)=\det(I+Z^{1/2}(P_0+R)Z^{1/2})^{-1/2}$，以区别于正文的下一时刻函数 $s(P)$。
每维分别使用 24 和 40 个节点，前者用于观察求积收敛，后者用于报告残差。
两工况共享 $P_0,R,Z,d$，因此该项积分检验对应共同的测量更新。
不同 $A$ 所产生的动态递推另由事件树和状态过程模拟检查。
上述积分仅检验条件高斯更新，不计算完整载荷日志的联合 TV。

\subsection{状态过程模拟}
另行模拟状态、测量及触发随机数，分别运行无攻击和受攻接收器，以核对事件树计算的损伤。
为降低配对损伤差的方差，每次实验的两个分支共享物理噪声及触发随机数，
并分别依据各自的实际预测计算发送决定。因此，两分支允许产生不同事件。
各次实验独立重置。取 $b=0.2$，每组进行 40 万次实验，使用预先固定的随机种子。
两拍策略的补偿动作发生在目标误差评价之后，因而不影响此处的 $J_k$。
结果见表\ref{tab:physical-validation}。
\begin{table}[H]
\centering\small
\tabcaption{状态过程模拟与事件树损伤的比较（$b=0.2$）}{State-process simulation and event-tree damage ($b=0.2$)}
\label{tab:physical-validation}
\begin{tabular}{llrrr}
\toprule
工况 & 策略 & 事件树 $J_k$ & 物理模拟均值 & 近似95\%置信区间\\
\midrule
S & 持续定向 & 0.0051947 & 0.0051157 & [0.0049781,0.0052534]\\
S & 两拍符号 & 0.0012874 & 0.0012858 & [0.0012289,0.0013427]\\
U & 持续定向 & 0.0125900 & 0.0127272 & [0.0124110,0.0130434]\\
U & 两拍符号 & 0.0022952 & 0.0023031 & [0.0022175,0.0023887]\\
\bottomrule
\end{tabular}\end{table}
表中区间由配对损伤差的样本均值加减 1.96 倍标准误得到，依赖正态近似；
附录\ref{app:statistical}的有限样本损伤保证则来自有界记录的 Hoeffding 不等式，二者适用对象不同。
实验次数、随机种子和标准误保存在 \texttt{vector\_data/physical\_check.json}。

\subsection{预算可行性与数值精度}
表\ref{tab:vector-constants}列出式\eqref{eq:attainable}及幅值上界所用的构造常数。
\begin{table}[H]
\centering\small
\tabcaption{二维算例中的解析构造常数}{Constants used in the two-dimensional attack constructions}
\label{tab:vector-constants}
\begin{tabular}{lrrrr}
\toprule
工况 & $L_v$ & $C_\pm$ & $C_{\rm peak}$ & $\kappa$\\
\midrule
S & 0.032185 & 3.663689 & 3.367727 & 1.000000 \\
U & 0.057380 & 3.511111 & 13.269454 & 1.000000 \\
\bottomrule
\end{tabular}
\end{table}

完整日志预算构造使用式\eqref{eq:Cpm}的显式 $C_\pm$，
乘以 $\eta^2$ 得到 TV 的充分上界，再与零动作策略进行私有随机混合以满足预算。
正文上下界图采用这些解析表达式的浮点计算结果，所示区间为最坏损伤的上下界，
并非对全策略最优值或连续载荷联合 TV 的精确数值求解。
二维表格采用普通浮点计算，不附带严格区间误差界。

\subsection{精确标签与记录界常数的一次精化}
\label{app:label-refinement}
将正文十进制参数视为精确有理数，按
$P^+=(P^{-1}+R_\gamma^{-1})^{-1}$ 及预测公式枚举，
S、U工况各32条历史均产生32个不同的后验协方差矩阵。
相等判定使用有理数运算，故精确标记与事件串的信息等价不依赖浮点去重阈值。

下面只评估标签依赖系数对现有记录界常数的影响。
对历史 $h=\gamma_{0:k}$，令 $M_t=(I-K_{\gamma_t}(P_t^-))A$，
$T_{j,h}=M_k\cdots M_{j+1}K_1(P_j^-)$（空乘积为 $I$）。偏差递推给出
$D_k=-\sum_{j:\gamma_j=1}T_{j,h}a_j$。于是取
\begin{equation}
 B(P)=b_{\rm ref}\max_{h:P_k^+(h)=P}\sum_{j:\gamma_j=1}\|T_{j,h}\|_2,
 \qquad \|D_k\|_2\le B(P_k^+).
 \label{eq:label-radius}
\end{equation}
此半径适用于全部允许策略。令 $h(P)=\lambda_{\min}(H(P))$，定义
\begin{equation}
 \begin{aligned}
 c(P)&=\frac{B(P)^2}{1-e^{-h(P)B(P)^2}}\quad(B(P)>0),\\
 F(P,Y)&=\operatorname{tr}P+c(P)[1-Y/s(P)].
 \end{aligned}
 \label{eq:label-score}
\end{equation}
$B(P)=0$ 时先取连续延拓 $c(P)=1/h(P)$。
割线界给出 $c(P)(1-e^{-D^\top H(P)D})\ge\|D\|_2^2$。
由条件静默恒等式，健康时 $\E_0F=\E_0\operatorname{tr}P$，因而
\begin{equation}
 J_k\le\E_\pi F-\E_0F\le C_{\rm label}\Delta_{\rm marked},
 \qquad C_{\rm label}=\max_{P,Y}F-\min_{P,Y}F.
 \label{eq:label-bound}
\end{equation}
最大、最小值取全部可达精确标签及 $Y\in\{0,1\}$。
若 $B(P)=0$，该标签下所有允许偏差均为零，亦可直接取 $c(P)=0$；
这是另一个有效取值，不是式\eqref{eq:label-score}的连续延拓。
表\ref{tab:label-constants}比较两种系数取值与原统一常数。
\begin{table}[H]
\centering\small
\tabcaption{标签依赖系数对记录界常数的影响}{Effect of label-dependent coefficients on the record-bound constant}
\label{tab:label-constants}
\begin{tabular}{lrrr}
\toprule
工况 & 原 $C_{\rm record}$ & 连续延拓 $C_{\rm label}$ & 零半径处取零\\
\midrule
S & 6.884171 & 6.710895 & 6.660168\\
U & 5.869119 & 5.303361 & 5.165567\\
\bottomrule
\end{tabular}
\end{table}
连续延拓取值使常数分别下降约2.52\%和9.64\%；零半径处取零时分别下降约3.25\%和11.99\%。
后一取值下，线性区间的上、下界斜率比仍约为758.14和316.09。
这些浮点评估说明，标签依赖系数仅有限改善当前常数，尚未消除上下界的明显间隙。
正文图\ref{fig:vector-budget}保留原常数，以对应主定理的统一构造。
精确标签表、两种边界取值及重建脚本见复现附件；常数小数值未作严格区间认证。


\subsection{复现文件与运行方式}
附件包含 Python 源代码、固定参数及完整的 CSV/JSON 结果，计算依赖 NumPy 和 Matplotlib。
运行 \texttt{python experiments.py} 可重建二维表格与图，
运行 \texttt{python physical\_check.py} 可复算状态过程模拟。
版本、随机种子及源文件哈希保存在结果文件中，环境配置与稿件编译方法见附件说明。
\texttt{review\_checks}目录提供确定性重放与标签核查脚本。
正式S、U算例的入口为\texttt{experiments.py}；\texttt{vector\_core.py}由其作为函数库调用。
正文图\ref{fig:vector-budget}的图例与解析界转折点由\texttt{plot\_review.py}从原参数和预算数据重建，无需再次模拟。
